\section{Improvement of Self-expression}
\label{sec:improve_se}

Numerous studies aim to improve the ability of LLMs to faithfully express their knowledge in responses. To provide a clearer understanding of current research, we classify the approaches into two main categories: training-free approaches, which include \emph{Prompting}, \emph{decoding-time intervention}, \emph{sampling and aggregation}, and \emph{post-generation revision}, and training-based approaches including \emph{self-aware fine-tuning} and \emph{Self-supervised fine-tuning}. Take an overview of these methods in Fig.~\ref{fig:improve_se}.

\begin{figure}[t]
\centering
\includegraphics[width=0.95\textwidth]{fig6_improve_self_expression.png}
\caption{Improvement of self-expression, encompassing both training-based and training-free approaches.}
\label{fig:improve_se}
\end{figure}

\subsection{Training-free Approaches}
\label{sec:se_training_free}

\paragraph{Prompting.} Pre-training enables LLMs to retain factual knowledge, but it is less effective at teaching them how to compose individual facts to answer complex questions (Press et al., 2023), leading to challenges in fully express their internal knowledge. One promising approach to address this issue is by using well-designed prompting strategies, as summarized in Table~\ref{tab:prompt_se}. Chain-of-thought prompting (CoT) (Wei et al., 2022) encourages LLMs to engage in a step-by-step reasoning process by providing few-shot demonstrations. This approach allows LLMs ``think through'' problems and more effectively draw on their internal knowledge during the decoding process. Later research explores zero-shot prompting by simply adding the phrase \emph{Let's think step by step} to the prompt (Kojima et al., 2022) to encourage step-by-step thinking.

Further studies focuses on structuring the generation process more explicitly. Zhou et al. (2023a) introduce least-to-most prompting, where the model first breaks a complex problem into smaller sub-questions, solves them one by one, and then combines the answers to tackle the original problem. Similarly, Self-ask (Press et al., 2023) prompts the LLM to decide when follow-up questions are needed and, if so, generates both the question and answer iteratively. Drawing from the cognitive skill of abstraction (Lachmy et al., 2022), step-back prompting (Zheng et al., 2024) allows LLMs to identify high-level abstractions, such as key concepts or principles, and use them to guide the generation process. In zero-shot settings, Wang et al. (2023c) propose plan-and-solve prompting, where the model first outlines a plan and then executes it step-by-step. Zhao et al. (2024) introduce fact-and-reflection, where the model first recalls relevant knowledge and then reflects on it to arrive at the final answer.

\begin{table}[t]
\centering
\caption{Overview of prompting methods for improving self-expression. The content within \{\} represents the response of LLMs.}
\label{tab:prompt_se}
\small
\begin{tabular}{p{3.2cm} p{9.5cm}}
\toprule
\textbf{Method} & \textbf{Prompt} \\
\midrule
Chain-of-Thought (Wei et al., 2022) & Q: Where is the highest mountain in the world? A: The highest mountain is Mount Everest, and it is located on the border of Nepal and Tibet (China). \\
\addlinespace
Zero-shot CoT (Kojima et al., 2022) & Let's think step by step. \\
\addlinespace
Least-to-Most (Zhou et al., 2023a) & Q: Where is the highest mountain in the world? Sub-Q1: \{Which is the highest mountain in the world?\} Sub-A1: \{Mount Everest\} Sub-Q2: \{Where is Mount Everest?\} Sub-A2: \{border of Nepal and Tibet (China)\} \\
\addlinespace
Self-ask (Press et al., 2023) & Q: Where is the highest mountain in the world? Are follow up questions needed: \{Yes\} Q1: \{Which is the highest mountain in the world?\} A1: \{Mount Everest\} Are follow up questions needed: \{Yes\} Q2: \{Where is Mount Everest?\} A2: \{Border of Nepal and Tibet (China)\} \\
\addlinespace
Step-back (Press et al., 2023) & Q: Which school did Kaiming He attend in November 2010? Stepback Q: \{What was Kaiming He's education history?\} Stepback A: \{B.S. in THU, 2007; Ph.D. in CUHK, 2011\} Final A: \{The Chinese University of Hong Kong\} \\
\addlinespace
Plan-and-solve (Wang et al., 2023c) & Let's first understand the problem and devise a plan to solve the problem. Then let's carry out the plan and solve the problem step-by-step. \\
\addlinespace
Fact-and-reflection (Zhao et al., 2024) & Q: Where is the highest mountain in the world? What are the facts needed to answer this question? Facts: \{1.The highest mountain is Mount Everest. 2.Mount Everest is located on the border of Nepal and Tibet (China).\} What is your reasoning? Reasoning: \{The highest mountain, Mount Everest, is located on the border of Nepal and Tibet (China).\} Final Answer: \{Border of Nepal and Tibet (China)\} \\
\bottomrule
\end{tabular}
\end{table}

\emph{Summary \& Discussion.} CoT prompting encourages LLMs to express their internal knowledge through step-by-step generation. However, its success depends largely on well-crafted prompts, which may be suboptimal and lack clear explainability. Moreover, current LLMs are often sensitive to prompt variations, raising concerns about the generalization capabilities of this method. Future research could explore the underlying mechanisms behind the effectiveness of CoT prompting or develop methods for automatically creating prompts that are both optimal and robust.

\paragraph{Decoding-time Intervention.} Numerous studies concentrate on eliciting the internal knowledge of LLMs through the decoding process. Li et al. (2024a) reveal a substantial gap between generation accuracy and probing accuracy as measured by a classifier on the hidden states of LLMs, suggesting that LLMs ``know'' more than they ``say''. To fully leverage this potential, they modify the activations of LLMs using truthful directions derived from the hidden states of true and false statements. Chen et al. (2024c) discover a strong correlation between the probability of hallucinations and contextual activations, which represents the mapping between the hidden states of context tokens and the predictive token. Using this information, they steer the decoding process to effectively reduce the occurrence of hallucinations.

Instead of investigating the hidden states, many studies directly modify the predictive distribution. Notably, contrastive decoding is extensively explored (Li et al., 2023a), where the logit difference between an expert model and an amateur model is utilized to steer generation, which amplifies the advantages of the expert and reduce the disadvantages of the amateur. Specifically, Chuang et al. (2024) employ the predictive distribution from the higher layer as the expert and that from the lower layer as the amateur, aiming to emphasize the factual knowledge embedded in the higher layer. Similarly, Zhang et al. (2023a) utilize the original LLM as the expert and a hallucination-prone LLM as the amateur to amplify the knowledge within the original model and reduce its tendency to fabricate information. Besides, Shi et al. (2024a) adopt the strong and weak activations of an MoE model as the expert model and the amateur model respectively, aiming to emphasize the reasoning capability of the MoE model's strong activations. In addition to focusing on parametric knowledge, Shi et al. (2024b) propose context-aware decoding, where both the expert and the amateur share the same LLM, but only the expert has access to context, thereby highlighting the importance of contextual knowledge in responses. The idea of contrastive decoding is also applicable in multimodal scenarios. Leng et al. (2024b) use a model with original visual inputs as the expert and the same model with distorted visual inputs as the amateur, with the goal of highlighting the role of visual inputs in shaping responses.

\emph{Summary \& Discussion.} Decoding-time intervention shows great promise in unlocking the potential of LLMs, but current research faces two primary challenges. First, the generalization is limited, as most methods have been concentrated on specific domains, such as factuality or reasoning, leaving their effectiveness in general instruction-following scenarios relatively underexplored. Additionally, these approaches may incur extra computational costs, such as an extra forward pass in contrastive decoding. Future research could develop strategies to address these challenges.

\paragraph{Sampling and Aggregation.} A straightforward approach for eliciting faithful knowledge from LLMs is through sampling and aggregation. This involves sampling multiple outputs and then aggregating them to derive the most consistent one, which is expected to more accurately reflect the model's knowledge. Wang et al. (2023d) sample a set of reasoning paths for a single query and then aggregate their answers by majority voting. This simple strategy achieves great performance in various reasoning tasks. Instead of aggregating answers, Chen et al. (2023b) prompt an LLM to aggregate multiple free-form responses based on majority consensus for open-ended generation. To make the most of multiple outputs, Thirukovalluru et al. (2024) split each output into several atomic parts, cluster them using sentence embeddings, remove clusters with fewer elements, and summarize the remaining clusters to produce a final consistent output. Similarly, Wang et al. (2024c) propose prompting an LLM to integrate and derive the final output based on the majority consensus from multiple outputs.

\emph{Summary \& Discussion.} Intuitively, consistent content across multiple generations tends to faithfully reflect the model's knowledge. While this approach demonstrates strong performance, the multiple sampling process incurs substantial computational costs, limiting its applicability in real-world settings. To overcome this challenge, future research could focus on constructing training data through sampling and aggregation, then fine-tuning the model with this data to internalize this capability.

\paragraph{Post-generation Revision.} Another approach involves post-generation refinement, where the response is modified to reduce inconsistencies with the model's knowledge. Dhuliawala et al. (2023a) first prompt the LLM to generate a list of questions designed to verify the atomic facts in its initial response, and then provide answers to these questions individually. Following this, the LLM is prompted to check the consistency between the initial response and the answers, making any necessary revisions. Similarly, Varshney et al. (2023) identify key concepts from the response, such as entities and keywords, evaluate their confidence using token probability, and retrieve external knowledge to validate and revise low-confidence concepts, which are more likely to be fabricated. Analogously, Zhao et al. (2023) assess the consistency across multiple responses as described by Wang et al. (2023d), then revise the less consistent ones using external knowledge.

\emph{Summary \& Discussion.} Post-generation revision offers an additional chance to correct inaccurate knowledge expression. The primary concern, however, is the increased computational cost. A possible direction for future research is to create training data based on this strategy and then fine-tune the LLM accordingly. For example, the original outputs and their revised versions could be treated as preference pairs, allowing methods like DPO or PPO to align the LLM with the desired attributes of the revised outputs.

\subsection{Training-based Approaches}
\label{sec:se_training_based}

\paragraph{Self-aware Fine-tuning.} Recent research suggests that fine-tuning LLMs with new knowledge may diminish their ability to express knowledge accurately, as this process can teach them to fabricate information beyond their internal knowledge (Gudibande et al., 2023; Lin et al., 2024a; Gekhman et al., 2024). To address this issue, many studies have started taking the knowledge boundaries of LLMs into account during fine-tuning, an approach we refer to as self-aware fine-tuning. Yang et al. (2024b); Zhang et al. (2024a); Cheng et al. (2024) fine-tune LLMs to explicitly state \emph{``I don't know''} when they lack adequate knowledge, thereby reducing the risk of generating fabricated information. Alternatively, Wan et al. (2024a) propose discard tuning, where examples are discarded when the model lacks the necessary knowledge, and open-book tuning, which incorporates reference knowledge into the context during fine-tuning to prevent the models from learning to fabricate content. More details on the knowledge identification methods used in these studies can be found in \S\ref{sec:sk_training_based}. Kang et al. (2024) introduce a conservative reward model that encourages less detailed responses in situations where the LLM is unfamiliar with the queries.

\emph{Summary \& Discussion.} Self-aware fine-tuning has demonstrated great potential in alleviating the tendency to fabricate information. The primary challenge lies in distinguishing between what the model knows and doesn't know. For future research, RL-based self-ware fine-tuning could be a valuable area for exploration, as it allows the LLM to explore broader spaces during training, potentially providing a more accurate reflection of its knowledge boundary.

\paragraph{Self-supervised Fine-tuning.} Another line of research fine-tunes LLMs to improve the expression of knowledge by leveraging supervision from their internal knowledge. Tian et al. (2024) initiate this process by prompting GPT-3.5 to extract multiple atomic claims from long-form generations. For each claim, they create a verification question, sample multiple answers from the LLM, and then calculate the consistency score across these answers. Based on the average consistency score of all claims, they construct preference pairs and employ DPO for optimization. Similarly, Zhang et al. (2024c) propose a method that use self-evaluation to verify each claim and also utilizes DPO for optimization. Simply, Lin et al. (2024a) prompt the LLM to generate responses for fact-based instructions, which are then used to fine-tune the model. After the fine-tuning phase, they create preference pairs by sampling responses from the LLM and use the model itself to assess these responses. Finally, they apply DPO for further refinement.

\emph{Summary \& Discussion.} Self-supervised fine-tuning has proven effective in improving LLMs' knowledge expression ability. However, a major concern is the quality of self-provided supervision signals, as these models may produce incorrect statements, and recent studies (Huang et al., 2024a; Jiang et al., 2024) have raised doubts about their self-evaluation abilities. Future research could aim to more thoroughly investigate the reliability of self-supervised fine-tuning.

\paragraph{Others.} In addition to the previously mentioned methods, there are other techniques aimed at improving the ability of LLMs to express their knowledge. Wei et al. (2023) construct simple synthetic data to train LLMs to avoid sycophancy, they reformat existing publicly available NLP datasets by adding user opinions that are independent of the correctness of the final answer and then fine-tune the models using this data. To achieve consistent model outputs across various prompts, Zhou et al. (2022); Cao et al. (2024) apply a consistency loss, which regularizes the outputs of semantically equivalent prompts to remain the same.
